\color{unidarkgreen}{\subsection{Логика отключения выключателя}}
\color{black}

\begin{enumerate}

\item
Функциональный блок логики отключения предназначен для подачи сигнала на отключение силового выключателя и формирования сигнала запрета АПВ силового выключателя. 

\item
В случае присоединения линии к системе через два выключателя для каждого из них можно задействовать свою логику отключения. Далее приведено описание логики отключения для выключателя В1, логика отключения для второго выключателя аналогична.

\item
Функциональная схема логики отключения В1 приведена на рисунке \ref{pict_lo:logic_otkl}. Функциональная схема логики отключения В2 выполняется аналогично схеме логики отключения В1. Уставки логики отключения В1 представлены в таблице \ref{app_set:lo_q1}. Уставки логики отключения В2 представлены в таблице \ref{app_set:lo_q2}.

\vspace{2mm}
\begin{figure}[H]
\centering
\includegraphics[width=1\textwidth,height=1\textheight,keepaspectratio]{\fbpath/03.01 ЛВ/761.9901 ЛОВ/img/logic_otkl.pdf}
\captionof{figure}{Функциональная схема логики отключения В1 и запрета АПВ В1}\label{pict_lo:logic_otkl}
\end{figure}

\item
Команда на отключение выключателя \textit{«Отключить»} формируется:
\begin{itemize}
    \item при срабатывании внутренних защит устройства (сигнал \textit{«Ср. защит»});
    \item при срабатывании ЗНФ (опционально, разрешается при введенном программном ключе \textbf{«ЗНФ В1 на откл.»});
    \item при действии УРОВ на отключение «своего» выключателя (сигнал \textit{«Срабатывание на себя»});
    \item от сигнала об аварийном давлении элегаза в ТТ, если введен программный ключ \textbf{«Откл. от Авар. давл. SF6 ТТ В1»};
    \item от команды оперативного отключения выключателя (сигнал \textit{«РКО»}).
\end{itemize}

\item
При необходимости может быть введен подхват команды на отключение от токовых реле, установленных в цепях ЭМО, в этом случае сигнал на отключение удерживается в сработанном состоянии до исчезновения тока через ЭМО (пока блок-контакты выключателя не разорвут цепи отключения). Ввод осуществляется программным ключом \textbf{«Подхват откл. В1 от РТ ЭМО»}.

\item
При наличии сигнала \textit{«Блок. отключения»}, формируемого функциональным блоком контроля силового выключателя, а также при выводе действия на выключатель отключение выключателя блокируется.

\item
Выдержка времени на возврат \textbf{«Тпрод откл. В1»} задает время продления сигнала отключения. 

\item
Команда на запрет АПВ выключателя \textit{«Запрет АПВ»} формируется безусловно:
\begin{itemize}
    \item от сигнала об аварийном давлении элегаза в ТТ (сигнал \textbf{«Откл. от Авар. давл. SF6 ТТ»});
    \item от команды оперативного отключения выключателя (сигнал \textit{«РКО»});
    \item от общей логики отключения.
\end{itemize}

\item
Команда на запрет АПВ может быть сформирована (при введеном соответствующем программном переключателе или выполнении определенных условий):
\begin{itemize}
    \item при срабатывании ЗНФ (сигнал \textit{«ЗНФ В1/ЗНФ: Срабатывание»});
    \item после фиксации отключенного состояния выключателя по окончании набора выдержки времени \textbf{«Тожид усл. пуска АПВ В1»};
    \item в режиме ведомого выключателя при включении от АПВ с контролем протекания тока через ЭМО.
\end{itemize}


\end{enumerate}